% Comandos para crear la portada y la primera página del PFG
%% Título del PFG
\newcommand{\PFGtitulo}{Planificación eficiente de infraestructuras virtuales basadas en docker y kubernetes para el streaming de vídeo}
%% Autor/a y directores/as del PFG
\newcommand{\PFGautor}{Ángel Roberto García Serpa}
\newcommand{\PFGdirector}{Agustín Carlos Caminero Herráez}
%% Descomentar la siguiente línea si hay codirector/a en el PFG
%\newcommand{\PFGcodirector}{Codirector/a del PFG}
%% Grado. A elegir entre:
%% Ingeniería Informática
%% Ingeniería en Tecnologías de la Información
\newcommand{\PFGgrado}{Ingeniería Informática}
%% Modalidad. A elegir entre:
%% genérica
%% específica
\newcommand{\PFGmodalidad}{genérica}
%% Curso académico en formato 20xx/20xx+1. Por ejemplo 2022/2023
\newcommand{\PFGcursoacademico}{2022/2023}
%% Convocatoria. A elegir entre:
%% junio
%% septiembre
%% diciembre
\newcommand{\PFGconvocatoria}{junio}
%% Fecha de lectura (si no se sabe el día concreto, poner mes y año}
\newcommand{\PFGfechadefensa}{Fecha de lectura}


\documentclass[12pt]{book}
\usepackage{graphicx}
\usepackage[utf8]{inputenc}
\usepackage[papersize={210mm,297mm},tmargin=25mm,bmargin=25mm,lmargin=20mm,rmargin=20mm]{geometry}

% no indentar los párrafos
\setlength\parindent{0cm}
\setlength\parskip{1em}

%Para tener los elementos de texto en español
\usepackage[spanish]{babel}
%Hipervinculos para todo lo útil
%\usepackage{hyperref}
\usepackage{xcolor}
\definecolor{links}{rgb}{0.35, 0.4, 0.35}
\definecolor{urls}{rgb}{0,0,1}
\definecolor{cites}{rgb}{0.8,0,0}

\usepackage{hyperref}
%\usepackage[all]{hypcap}
\hypersetup{
	colorlinks = true,
	linkcolor = links,
	urlcolor = urls,
	citecolor = cites
}

\usepackage{fancyhdr}

%añadir ficheros externos
\usepackage{subfiles}
%soporte para anexos
\usepackage{appendix}
%paquetes para figuras
\usepackage{caption}
\usepackage{subcaption}
\usepackage{listings}

%Encabezados bonitos
\pagestyle{headings}

%paquete para poner la bibliografia en la table of contents
\usepackage[nottoc,notlot,notlof]{tocbibind}

%Cambiamos el espacio en blanco enorme que por defecto hay en las páginas inicio de capítulo
\usepackage{titlesec}
\titleformat{\chapter}[display]   
{\normalfont\huge\bfseries}{\chaptertitlename\ \thechapter}{20pt}{\Huge}   
\titlespacing*{\chapter}{0pt}{-50pt}{40pt}

%para crear subsubsubsection
\setcounter{secnumdepth}{4}
\titleformat{\paragraph}
{\normalfont\normalsize\bfseries}{\theparagraph}{1em}{}
\titlespacing*{\paragraph}
{0pt}{3.25ex plus 1ex minus .2ex}{1.5ex plus .2ex}

%Ruta en la que pondremos las imágenes para el documento
\graphicspath{ {imagenes/} }
%Redifinimos el nombre que asigna babel spanish a las tablas (cuadro) a tabla
\renewcommand{\spanishtablename}{Tabla}
\renewcommand{\spanishlisttablename}{Índice de tablas}
\renewcommand{\spanishcontentsname}{Índice}
\renewcommand{\appendixname}{Anexos}
\renewcommand{\appendixtocname}{Anexos}
\renewcommand{\appendixpagename}{Anexos}

% Definimos el comando de página en blanco
\newcommand\paginablanco{%
    \null
    \thispagestyle{empty}%
    \newpage}

% Definimos el comando de página en blanco sin avanzar numeracion
\newcommand\paginablancosin{%
    \paginablanco{}
    \addtocounter{page}{-1}}

%mayor anchura en las tablas
\renewcommand{\arraystretch}{1.5}
%mayor ancuhura en las fracciones
\newcommand\ddfrac[2]{\frac{\displaystyle #1}{\displaystyle #2}}
%comando escribir TFG
\newcommand{\tfg}{Trabajo Fin de Grado }
%definimos estilo consola
\definecolor{mygreen}{rgb}{0,0.6,0}
\definecolor{mygray}{rgb}{0.5,0.5,0.5}
\definecolor{mymauve}{rgb}{0.58,0,0.82}
\definecolor{terminalbgcolor}{HTML}{330033}
\definecolor{terminalrulecolor}{HTML}{000099}
\newcommand{\lstconsolestyle}{
	\lstset{
		backgroundcolor=\color{terminalbgcolor},
		basicstyle=\color{white}\fontfamily{fvm}\footnotesize\selectfont,
		breakatwhitespace=false,  
		breaklines=true,
		captionpos=b,
		commentstyle=\color{mygreen},
		deletekeywords={...},
		escapeinside={\%*}{*)},
		extendedchars=true,
		frame=single,
		keepspaces=true,
		keywordstyle=\color{blue},
		%language=none,
		morekeywords={*,...},
		numbers=none,
		numbersep=5pt,
		framerule=2pt,
		numberstyle=\color{mygray}\tiny\selectfont,
		rulecolor=\color{terminalrulecolor},
		showspaces=false,
		showstringspaces=false,
		showtabs=false,
		stepnumber=2,
		stringstyle=\color{mymauve},
		tabsize=2
	}
}


% Comenzamos el documento
\begin{document}
\pagenumbering{gobble} % para no incluir números en la portada y primera página

% Hacemos la portada
\begin{titlepage}

\begin{figure}[!htb]
    \centering
    \includegraphics{rosadelosvientos.png}
\end{figure}

\begin{center}
    {\scshape\Large Universidad Nacional de Educación a Distancia \par}
    {\scshape\Large Escuela Técnica Superior de Ingeniería Informática \par}
    \vfill
    \large{Proyecto de fin de Grado en \PFGgrado}\\
    \vfill
    {\bfseries\Huge \textcolor{black}{\PFGtitulo} \par}
    \vfill
\end{center}

\vspace*{3cm}
{\huge \textcolor{black}{\PFGautor} \par}
{\huge \textcolor{black}{Dirigido por: \PFGdirector} \par}
\ifdefined\PFGcodirector
  \huge \textcolor{black}{Codirigido por: \PFGcodirector} \par
\fi
{\huge \textcolor{black}{Curso \PFGcursoacademico, convocatoria \PFGconvocatoria} \par}

\chapter*{}

\thispagestyle{empty}

\begin{figure}[!htb]
    \centering
    \includegraphics{rosadelosvientos.png}
\end{figure}

\begin{center}

    {\bfseries\Huge \textcolor{black}{\PFGtitulo} \par}
    \vfill
    {\bfseries\Large \textcolor{black}{Proyecto de fin de Grado en \PFGgrado} \par}
    {\bfseries\Large \textcolor{black}{de modalidad \PFGmodalidad} \par}
\end{center}

\vspace*{5cm}
{\Large Realizado por: \PFGautor \par}
{\Large Dirigido por: \PFGdirector \par}
\ifdefined\PFGcodirector
  {\Large Codirigido por: \PFGcodirector \par}
\fi%

%%{\Large Codirigido por: \PFGcodirector \par}
\vfill
{\Large Fecha de lectura y defensa: \PFGfechadefensa \par}

\end{titlepage}

\pagenumbering{Roman} % para comenzar la numeracion de paginas en numeros romanos

% Agradecimientos
\chapter*{Agradecimientos}
\thispagestyle{empty}
\begin{flushright}
\textit{Incluir todos los agradecimientos}
\end{flushright}

\chapter*{Resumen} % El * indica que no queremos que añada la palabra "Capitulo"

Se pretende realizar despliegue y evaluación de diversas tecnologías de servidores de streaming de video sobre contenedores ligeros. Para la consecución del objetivo principal, se deberán realizar los siguientes:
\begin{list}{·}{}
	\item Creación de las infraestructuras paralelas distribuidas.
	\item Recopilación de información sobre servidores de streaming de video.
	\item Puesta en marcha y evaluación de prestaciones de los servidores de vídeo seleccionados.
	\item Generación de documentación para el despliegue de los servidores utilizados tanto en local como en la nube.
\end{list}

\chapter*{Abstract} % El * indica que no queremos que añada la palabra "Capitulo"

Resúmen del PFG en inglés.

\chapter*{Palabras clave} % El * indica que no queremos que añada la palabra "Capitulo"

video, distribuido, Kubernetes 

\chapter*{Keywords} % El * indica que no queremos que añada la palabra "Capitulo"

video, distributed, Kubernetes

\tableofcontents %% Índice general
\listoftables    %% Índice de tablas
\addcontentsline{toc}{section}{Índice de tablas} % Queremos que aparezca en el índice
\listoffigures   %% Índice de figuras
\addcontentsline{toc}{section}{Índice de figuras} % Queremos que aparezca en el índice
\newpage

\paginablanco{}

%% Comenzamos con la numeración arábica porque aquí empieza el contenido
\pagenumbering{arabic}

%% Ejemplo de importación de capítulos
\chapter{Introducción}
\subfile{capitulos/C1-introduccion}

\chapter{Estado del arte}
\subfile{capitulos/C2-estado_arte}
\chapter{Propuesta}
\subfile{capitulos/C3-propuesta}
\chapter{Diseño}
\subfile{capitulos/C4-disenno}
\chapter{Desarrollo del proyecto}
\subfile{capitulos/C5-desarrollo}
\chapter{Pruebas}
\subfile{capitulos/C6-pruebas}
\chapter{Conclusiones}
\subfile{capitulos/C7-conclusiones}

%% Bibliografía
\newpage
\subfile{capitulos/referencias}

\clearpage

\appendix
\chapter{Primer anexo}
\subfile{capitulos/A1-anexo}

\chapter{Segundo anexo}
\subfile{capitulos/A2-anexo}

%...

\end{document}